\documentclass[11pt]{article}
\usepackage[a4paper,margin=24mm]{geometry}
\usepackage{amsmath,amssymb,amsthm,booktabs,hyperref}
\newtheorem{theorem}{Theorem}\newtheorem{lemma}[theorem]{Lemma}
\newcommand{\F}{\mathbb F}\newcommand{\Om}{\Omega}\newcommand{\PGU}{\operatorname{PGU}}\newcommand{\GL}{\operatorname{GL}}\newcommand{\End}{\operatorname{End}}\newcommand{\PSU}{\operatorname{PSU}}
\title{Maximal overgroups of a full negative $6+2$ torus element}
\author{L045 recovery: a classification interface}\date{}
\begin{document}\maketitle
Let $q=p^f\ge4$, $d=(2,q-1)$, $A=q^3+1$, and $B=q+1$.
Let $V=V_6^-\perp V_2^-$ over $\F_q$, and let
$X=g\perp h$ where $g$ and $h$ generate the norm-one tori of orders
$A$ and $B$ on the two summands. Thus $X\in\Om_8^+(q)$ and its image
$x\in G=P\Om_8^+(q)$ has order $A/d$, by the acquired torus proof.
The primitive elements $g,h$ are arbitrary; no relation between their
choices is assumed.

\begin{theorem}\label{thm:cover}
Every maximal subgroup of $G$ containing $x$ is one of the following:
\begin{enumerate}
\item the unique stabilizer $D$ of $V_6^-\perp V_2^-$;
\item the normalizer of a unitary structure on a four-dimensional
$\F_{q^2}$-space; there are at most two such subgroups containing $x$;
\item a subgroup of type $B_3$: $\Om_7(q)$ for odd $q$, or
$\operatorname{Sp}_6(q)$ for even $q$, in one of the three triality actions.
\end{enumerate}
The natural-action $B_3$ subgroups do not contain $x$.
Also $N_G(\langle x\rangle)\le D$, so a separate cyclic-normalizer
term is not missing from this cover.
\end{theorem}

The only classification premise is the complete published maximal-subgroup
classification in Bray--Holt--Roney-Dougal (BHRD), Main Theorem 2.1.1,
Section 8.1 and Table 8.50 \cite{BHRD}. The table is explicitly attributed
there to Kleidman's classification of these orthogonal groups.
We use the table for the simple group $G$, not for arbitrary almost simple
extensions. Its novelty rows therefore give no additional maximal subgroup
of $G$. All eliminations for our particular elements are proved below.
The complete row-by-row routing is retained separately.

\begin{lemma}\label{lem:primitive}
There is a prime $r\mid |x|$ with $\operatorname{ord}_r(p)=6f$.
Every such prime satisfies $r\ge6f+1$ and $\operatorname{ord}_r(q)=6$.
Moreover the entire primitive part $\Phi_{6f}^*(p)$ divides $|x|$.
\end{lemma}
\begin{proof}
The primitive-divisor theorem applied to $p^{6f}-1$ has no exception here:
the exceptional pair $(p,6f)=(2,6)$ would mean $q=2$.
All primitive prime powers divide $p^{3f}+1=A$, and none divides $2$.
This proves the assertions after passing to the projective quotient of order
$d\le2$. We retain the classical primitive-divisor theorem as a named
number-theoretic input. Its exact statement and the primitive-part convention
are also recalled in the acquired Bamberg--Penttila preprint
\cite[pp.2--3]{BP}.
\end{proof}

\begin{proof}[Proof of the type assertion in Theorem~\ref{thm:cover}]
Apply the complete BHRD table to the preimage of a maximal subgroup in
$\Om_8^+(q)$. Central factors of order at most two do not affect the
primitive divisor. The parabolic rows have Levi factors of types
$\GL_4(q)$, $\GL_2(q)\Om_4^+(q)$ or $\GL_3(q)$; their orders have no
factor $r$. The same test excludes the positive $2+6$ and $\operatorname{SL}_4$
rows, the $3+5$ and $\operatorname{Sp}_2\circ\operatorname{Sp}_4$ rows,
the four-plane decompositions, the two four-space decompositions, and
$\Om_4^+(q^2)$. Their non-characteristic order factors divide products of
$q^j-1$ for $j\le4$, together with powers of $2$ and $3$.
The additional dihedral novelty row has factors $q^2+1$ and powers of two,
and is excluded by the same argument.
The negative $2+6$ and unitary four-space rows survive.

The subfield rows are excluded using the \emph{absolute} order $6f$.
For $q=q_0^s$, $q_0=p^a$, $s\ge2$, the factors in $\Om_8^+(q_0)$
have exponents at most $6a<6f$.
For the two $\Om_8^-(q_0)$ actions with $q=q_0^2$, the exponents are
at most $8a<12a=6f$.
For ${}^3D_4(q_0)$ with $q=q_0^3$, use
\[
 |{}^3D_4(q_0)|=q_0^{12}(q_0^8+q_0^4+1)(q_0^6-1)(q_0^2-1),
 \qquad q_0^8+q_0^4+1\mid q_0^{12}-1.
\]
Here $12a<18a=6f$. All displayed extension factors in these table rows
have order a power of two. The $L_3(q).3$ row has factors of exponents
at most three and an extra factor three, so is excluded too.

Consider the $U_3(q).3$ row, which occurs only when $q\equiv2\pmod3$.
In its projective group, $y=x^3$ lies in $\PSU_3(q)$ and has order
$A/(3d)$. An $r$-element in $\PSU_3(q)$ acts irreducibly on its natural
three-space over $\F_{q^2}$, since $\operatorname{ord}_r(q^2)=3$.
Its centralizer is the norm-one Singer torus modulo scalars, of order
$(q^2-q+1)/3$. This follows directly by identifying the commuting
algebra with $\F_{q^6}$: the unitary condition is norm one to
$\F_{q^3}$, the determinant-one condition leaves order $q^2-q+1$,
and the scalar center has order three. Projective centralization gives no
extra elements: lift the $r$-element with order $r$; a commuting lift
can differ by a scalar of order three, which its $r$th power forces to be one.
Since $y$ centralizes its own $r$-subgroup, its order must divide this torus
order. This is impossible, their ratio being $B/d>1$.

The remaining bounded-order rows are handled as follows. The monomial and
extraspecial rows have a normal $2$-group with quotient contained in $S_8$
(or $L_3(2)$), and require $q=p$ odd. Their primitive part is therefore at
most seven: a primitive divisor is at least seven, and seven occurs only
once in those quotient orders. The same assertion holds for the row
$2\cdot\Om_8^+(2)$, whose prime divisors are $2,3,5,7$, with only one
factor seven. But
\[
 \Phi_6^*(p)=\frac{p^2-p+1}{(3,p+1)}>7\qquad(p\ge7).
\]
Only $q=5$ remains in our range. Now $|x|=63$. An odd-order element
injects through a $2$-group kernel; $S_8$ and $L_3(2)$ have no element of
order $63$. Neither does $\Om_8^+(2)$, as proved in
Lemma~\ref{lem:63} below. The exceptional rows $2\cdot\mathrm{Sz}(8)$
and the two central extensions of $A_{10}$ occur only at $q=5$.
The first group has order prime to three; the second ones cannot have
order $63$, since a permutation of that order requires disjoint cycles
of lengths nine and seven. The $A_9$ row has $q=2$ and is outside the
hypothesis. The $G_2$ row is marked as a novelty and is not maximal
in the simple group $G$.
Thus only the three stated types remain. A natural $B_3$ subgroup fixes
a nonsingular line, whereas $X$ has no eigenvalue in $\F_q$.
\end{proof}

\begin{lemma}\label{lem:63}
No element of the orthogonal group $O_8^+(2)$ has order $63$.
\end{lemma}
\begin{proof}
Such an element would be semisimple. An eigenvalue whose order is divisible
by $63$ has degree at least six, and an eigenvalue of order $63$ has degree
six. Its irreducible polynomial is not self-reciprocal: no power of $2$
is $-1$ modulo $63$. A second, distinct dual constituent would require
twelve dimensions. If no eigenvalue has order divisible by $63$, separate
constituents must supply the factors nine and seven. Order nine requires
six dimensions, since $\operatorname{ord}_9(2)=6$. A constituent supplying
the factor seven costs at least six dimensions: if its degree is three,
its eigenvalue has order seven and a distinct dual constituent is required;
otherwise its degree is at least six (this includes eigenvalue order $21$).
The two contributions are distinct, since otherwise an eigenvalue would
have order divisible by $63$. Again at least twelve dimensions are necessary.
\end{proof}

\begin{proof}[Proof of the multiplicities and normalizer assertion]
Both $X$ and $X^2$ have exactly two distinct irreducible constituents, of
dimensions six and two. Indeed $g^2$ has order $A/(2,A)$, which exceeds
$q^2-1$ and does not divide $q^3-1$; its degree is therefore six.
The order of $h^2$ is $B/(2,B)>2$ for $q\ge4$, so its degree is two.
Thus the only proper nonzero $X$-invariant subspaces are $V_6$ and $V_2$,
and the decomposition stabilizer $D$ is unique.

Let $E\cong\F_{q^2}$ be a quadratic scalar-field structure whose unitary
normalizer contains $X$. Conjugation by $X$ induces an automorphism of
$E$ of order at most two, so $X^2$ commutes with $E$. Consequently
\[
 E\ \subseteq\ C_{\End_{\F_q}(V)}(X^2)
       =\F_{q^6}\oplus\F_{q^2}.
\]
This algebra is commutative and contains $X$, so $X$ in fact centralizes
$E$. Each coordinate projection of $E$ is an embedding. Up to a common
Frobenius automorphism of $E$, there are exactly two choices, determined
by the relative Frobenius on the two factors. A field structure and the
quadratic form determine its full unitary extension normalizer; hence
at most two of the unitary maximal subgroups can contain $x$.

Finally a lift of an element normalizing $\langle x\rangle$ conjugates
$X$ to $\epsilon X^j$, where $\epsilon\in\{1,-1\}$ and
$(j,A/d)=1$. The two constituents of $X^j$ still have degrees six and two:
$(j,A)\le d$, so their orders lose at most a factor two, and multiplication
by $\epsilon$ does not change their degrees. The preceding degree argument
applies. Their unique dimensions force the lift to preserve both summands.
Therefore $N_G(\langle x\rangle)\le D$.
\end{proof}

\paragraph{Scope and attribution.}
This is an application of an explicitly retained complete subgroup
classification, not a new classification. It corrects the missing arbitrary-element
interface in the use of Leemans--Liebeck Table 4; it does not infer that
that source's multiplicities or projective torus orders are correct.
The triality analysis deciding which $B_3$ subgroups contain a particular
$x$, the inverter bounds and the final generating-pair inequality are
separate components. No whole-original admission, measurement or new
question is asserted here. The torus construction and the surrounding
orthogonal program were authored by the root agent; this cover specialization
and its two-field-structure argument were authored by publication\_recovery.
The parent independently obtained the relative-field orientation argument
in a separate orthogonal manuscript. The peer f003\_resume supplied a
supporting elementary order-$21/63$ note and arithmetic checks; these are
retained as a contribution, not as an independent review of this component.

\begin{thebibliography}{9}
\bibitem{BHRD} J. N. Bray, D. F. Holt and C. M. Roney-Dougal,
\emph{The Maximal Subgroups of the Low-Dimensional Finite Classical Groups},
LMS Lecture Note Series 407, Cambridge University Press, 2013.
Main Theorem 2.1.1, Section 8.1, Table 8.50 (printed pp.55,373--374,402--404).
Complete published book acquired from the National Academic Digital Library;
the classification is retained, not recomputed.
\bibitem{BP} J. Bamberg and T. Penttila,
\emph{Overgroups of cyclic Sylow subgroups of linear groups},
Comm. Algebra 36 (2008), 2503--2543; complete 27-page author-hosted preprint
acquired from Bamberg's Ghent preprint page, with its old-draft qualification
and current author errata explicitly retained. Only its elementary
primitive-part convention is used above; the complete maximal cover comes
from BHRD.
\end{thebibliography}
\end{document}
